\subsection{Some inseparable product polarizations: sharp rectangular ranks}
\label{sec:ss-rectangular-ranks}

\begin{lemma}[Cartier bounds from one or two even factors]
\label{lem:ss-rectangular-even-bound}
\status{Proved; independent full proof audit passed.}
Let $E:y^2+y=x^3$ over an algebraically closed field of characteristic
two, let $A=E\times E$, and put
$N=\OO_E(mO)\boxtimes\OO_E(nO)$, where $m,n\ge3$.
On $W=H^0(A,N)$ define the root-removed Cartier matrix $C_s$ by
the Cartier trace of $su\,dx_1\wedge dx_2$, for $s,u\in W$.
The generic stable rank on $H^2(A,N^{-1})$ is at most
$mn-2$ if $m$ is even, and at most $mn-3$ if both $m,n$ are even.
For every smooth divisor $Y\in|N|$, the stable rank on its relative
coherent cohomology equals the stable rank on $H^2(A,N^{-1})$.
\end{lemma}
\begin{proof}
The invariant derivation on $E$ is
$\partial x=1$, $\partial y=x^2$, and $\partial^2=0$.
If $m$ is even, the line bundle $\OO_E(mO)$ is pulled back by
relative Frobenius from $\OO_{E^{(2)}}((m/2)O)$.  Its canonical
connection therefore makes $\partial_1$ a global $k$-linear
endomorphism of $W$.  In horizontal local trivializations the
following identities are Cartier identities of rational differential
forms and hence identities of regular twisted sections:
\[
                  C_s(s)=0,\qquad C_s(\partial_1s)=0.
\]
For the first, $C(s^2\,dx_1\wedge dx_2)=sC(dx_1\wedge dx_2)=0$.
For the second, write $s=f$ locally.  Then
$f\partial_1f\,dx_1\wedge dx_2=f\,df\wedge dx_2$.
Both factors on the right are closed, so multiplicativity of Cartier
and $C(dx_2)=0$ give its vanishing.  The horizontal transition in
the first factor and the transition depending only on the second
coordinate make $\partial_1s$ a section of $N$; this is why parity
is required for this particular global relation.

Generically $s$ and $\partial_1s$ are linearly independent.
Indeed, if $\partial_1s=cs$ with $c\in k$, applying $\partial_1$
gives $c^2s=0$, hence $c=0$.  The operator is nonzero: the rational
section $x_1$ belongs to $W$ and has nonzero derivative, using the
horizontal rational trivialization with pole divisor $mO$.
If both factors are even, the same argument gives
$C_s(\partial_2s)=0$.  The three sections
$s,\partial_1s,\partial_2s$ are generically independent, as one sees
already from $s=x_1x_2$: they are $x_1x_2,x_2,x_1$.
The corresponding matrix-minor vanishing holds on a dense open and
therefore everywhere.  Linear rank bounds stable rank, proving both
inequalities.

The divisor sequence and ample vanishing give a Frobenius-compatible
exact sequence
\[
0\longrightarrow H^1(A,\OO_A)\longrightarrow H^1(Y,\OO_Y)
 \longrightarrow H^2(A,N^{-1})\longrightarrow H^2(A,\OO_A)
 \longrightarrow0.
\]
Frobenius on both ambient terms is zero, because $E$ is supersingular.
Passing to stable parts therefore identifies the stable part of the
full curve cohomology, of its relative quotient, and of
$H^2(A,N^{-1})$.  The latter Frobenius is Serre dual to the indicated
Cartier matrix, with the usual Frobenius twist.  These statements also
hold for reduced divisors used below: the same exact divisor sequence
and the same vanishing apply.
\end{proof}

\begin{lemma}[An exact rectangular base and its odd norm images]
\label{lem:ss-rectangular-base-norm}
\status{Proved; independent full proof audit and independent exact matrix reconstruction passed.}
For $N_0=\OO_E(6O)\boxtimes\OO_E(3O)$, a general smooth divisor
has genus $19$ and full stable rank $16$.
If $a,b$ are positive odd integers, a general smooth divisor of
$\OO_E(6aO)\boxtimes\OO_E(3bO)$ has full stable rank $18ab-2$.
\end{lemma}
\begin{proof}
We first supply a finite exact certificate for the base.  Use the bases
$(1,x,x^2,x^3,y,xy)$ and $(1,x,y)$, with tensor index $3i+j$.
Over $\mathbf F_2$ take the coefficient vector encoded by the integer
$74594$, its $j$-th coefficient being its $j$-th binary digit.
The eighteen columns of the Cartier matrix, encoded in the same way,
are
\[
\begin{gathered}
4106,12,16432,4177,40,16660,32848,96,131456,\\
33416,320,133280,12424,49472,82080,65664,164608,66560.
\end{gathered}
\]
Binary Gaussian elimination gives
$\operatorname{rank}C_s=\operatorname{rank}C_s^2=16$.
The exact code, both sampled inputs (including the rejected one),
complete matrix and output are
\path{code/supersingular_rectangular_rank.py} and
\path{outputs/supersingular_rectangular_rank.json}.
The formula used to build every column is
\[
 C(x^r dx)=\begin{cases}0&r\text{ even},\\x^{(r-1)/2}dx&r\text{ odd},\end{cases}
\quad
 C(x^r y\,dx)=\begin{cases}x^{r/2+1}dx&r\text{ even},\\
                       x^{(r-1)/2}y\,dx&r\text{ odd},\end{cases}
\]
with $y^2=y+x^3$ and the tensor product on the surface.
Thus this is a reproducible certificate, not an inference from
sample frequency.  Equality of the first two ranks makes the image
stable and proves stable rank $16$ at this coefficient vector.
The corresponding nonzero minors of iterated Frobenius define a
nonempty open.  The upper bound in
Lemma~\ref{lem:ss-rectangular-even-bound} is $16$, and $N_0$ is very
ample.  Bertini's nonempty smooth open intersects this rank open.
Adjunction gives genus $1+\chi(N_0)=19$.

For each positive odd integer $c$ there is an \'etale endomorphism of
$E$ of degree $c$.  Indeed Lemma~\ref{lem:DE-A3-ss-norm} constructs
an endomorphism of every positive degree by explicit quaternion norms
and the four-square identity.  An endomorphism of odd degree in
characteristic two is automatically \'etale.
Choose such endomorphisms of degrees $b,a$ on the first and second
factors, respectively, and let $\phi:A\to A$ be their product.
It is \'etale of degree $ab$.  Pushforward of a horizontal or vertical
divisor gives
\[
     \operatorname{Nm}_\phi(N_0)\equiv
                  \OO_E(6aO)\boxtimes\OO_E(3bO)=N_{a,b}.
\]

We prove that a general source curve has only nodes after applying
$\phi$, including kernels whose projections are not injective.
For a nonzero kernel element $g$, impose the two values
$s(P)=s(P+g)=0$.  They are independent because $N_0$ is very ample.
If $g_1\ne0$, degree six on the first elliptic factor separates the
two first jets, and degree three on the second separates a first
jet at one point.  Tensor products therefore give all six surface
first-jet coordinates at $P,P+g$.  Tangency imposes one further
nonzero determinant equation after the two value equations.
If $g_1=0$, then $g_2\ne0$.  Degree three on the second factor has
rank three on the length-four scheme $2P_2+2(P_2+g_2)$: its kernel
is zero because a nonzero degree-three section cannot vanish to
total order four.  Its two value functionals are independent.
The two derivatives in the first factor are independently supplied
by the two second-factor value functionals.  After imposing the
values, the two second-direction derivatives span a nonzero line;
otherwise a degree-three section would again vanish to order four.
The tangent determinant is consequently a nonzero polynomial in
these derivative parameters.  In either case, the incidence of
tangencies has dimension at most $\dim|N_0|-1$ after allowing the
two-dimensional point $P$ to vary.  There are finitely many $g$.

For two distinct nonzero kernel elements $g,h$, the three value
conditions at $P,P+g,P+h$ are independent unless all three first
coordinates coincide and their three second coordinates form a
degree-three divisor linearly equivalent to $3O$.
To verify this, separate the at most three distinct first coordinates
using the degree-six bundle, then separate the at most two second
coordinates within each repeated first-coordinate class using the
degree-three bundle.  Only the class of size three remains; on an
elliptic curve its value rank drops from three to two precisely when
the sum of its three points is zero.
That exceptional condition is $3P_2+g_2+h_2=0$, so its point locus
has dimension one, because $[3]$ is finite \'etale.  The ordinary
rank-three incidence over a dimension-two locus and the exceptional
rank-two incidence over a dimension-one locus both have dimension
at most $\dim|N_0|-1$.  Thus a general source member has no triple
identifications.  A higher identification contains a triple and is
also excluded.  The tangency calculation excludes a translated
component equal to the source, hence $\phi|_{C'}$ is birational onto
its image.  Its image is integral and nodal, because $\phi$ is \'etale
and all remaining pairs are transverse.

Choose $C'$ also in the smooth rank-$16$ open.  The image has
arithmetic genus $18ab+1$, normalization genus $19$, and exactly
$18(ab-1)$ nodes.  The normalization sequence contributes a
Frobenius-bijective vector space of that dimension.  Its full stable
rank is therefore $16+18(ab-1)=18ab-2$.
The two ample line bundles in the numerical equality above differ
by an element of $\Pic^0(A)$.  Surjectivity of the polarization
$\lambda_{N_{a,b}}:A\to A^\vee$ gives a target translation putting
the nodal image in the exact system $|N_{a,b}|$.  No separability of
this polarization, and no two-torsion translation, is needed for
this step.  Translation preserves every rank and singularity used.
Openness of iterated Frobenius rank and the nonempty smooth-member
open give the same lower bound for a general smooth divisor.
Lemma~\ref{lem:ss-rectangular-even-bound} gives the matching upper
bound, proving the result.
\end{proof}

\begin{theorem}[All degrees for some even product types]
\label{thm:ss-rectangular-all-degrees}
\status{Proved; independent full proof audit passed.}
Let $E:y^2+y=x^3$ over a perfect field of characteristic two, and
let $A=E\times E$, and let $a,b$ be positive integers.  Put
$M=\OO_E(aO)\boxtimes\OO_E(bO)$ and $L=M^3$.
The relative coherent dimension in $|L^d|$ is $9abd^2-1$.
If exactly one of $a,b$ is even and that even integer has
$2$-adic valuation one, then for every $d\ge1$ the geometric generic
relative stable rank is
\[
 r_d=\begin{cases}9abd^2-2&d\text{ odd},\\
                  9abd^2-3&d\text{ even}.
      \end{cases}
\]
If both $a,b$ are even, it is $r_d=9abd^2-3$ for every $d\ge1$.
For the embedding by $L$, these formulas also hold for ambient
hypersurfaces whenever
$d\ge\max\{1,\operatorname{reg}(\mathcal I_A)-1\}$.
\end{theorem}
\begin{proof}
Pass to an algebraic closure; geometric generic ranks are unchanged.
In the first case interchange the factors if necessary and write
$a=2a'$ with $a',b$ odd.  For odd $d$, apply
Lemma~\ref{lem:ss-rectangular-base-norm} with parameters $a'd,bd$.
It gives $18a'bd^2-2=9abd^2-2$.

Whenever $ad$ and $bd$ are both even, decompose the target divisor
class as
\[
 (3ad,3bd)=(3,3)+(3(ad-1),3(bd-1)).
\]
Both summands are cubes of ample product bundles with odd Euler
characteristic.  Theorem~\ref{thm:DE-A3-separable-ranks} supplies
smooth members in each system with full curve defect two.
Choose them simultaneously in their sharp opens and transverse
to each other; very ampleness and Bertini make these conditions
nonempty opens.  Their union is connected, since their intersection
number is positive.  If their genera are $g_1,g_2$ and they meet in
$q$ transverse points, normalization gives genus $g_1+g_2+q-1$
and stable rank $(g_1-2)+(g_2-2)+q-1$.  Thus the union has full
defect four and stable rank $9abd^2-3$.
The rank open again meets the smooth-member open in the target
linear system.  The both-even bound of
Lemma~\ref{lem:ss-rectangular-even-bound} is exactly $9abd^2-3$,
so equality holds.  This proves all even $d$ in the first case and
every $d$ in the second.

The relative dimension follows from surface Riemann--Roch and ample
vanishing.  Finally the indicated regularity bound makes ambient
restriction onto $H^0(A,L^d)$ surjective; the induced parameter map
has integral geometric fibres, preserving its geometric generic
rank.  The remaining product types with one even factor of higher
$2$-adic valuation are settled by
Theorem~\ref{thm:DE-A3-all-ss-rectangles} below. Arbitrary
inseparable polarizations on a nonproduct surface remain outside
these product-polarization results.
\end{proof}
